\documentclass[11pt]{article}
\usepackage[T1]{fontenc}
\usepackage{lmodern,amsmath,amssymb,microtype}
\usepackage[margin=1.02in]{geometry}
\usepackage[colorlinks=true,urlcolor=blue,citecolor=blue]{hyperref}
\setlength{\emergencystretch}{2em}
\hypersetup{pdfauthor={Wentao Li},pdftitle={Projective type and self-surjections of power sets}}
\title{Projective type and self-surjections of power sets}
\author{Wentao Li\\[2pt]\small Prepared with assistance from OpenAI Codex}
\date{November 2026}
\begin{document}
\maketitle
\begin{abstract}
Let $A$ be amorphous. In ZF, a positive finite power of $\mathcal P(A)$ admits a noninjective self-surjection if and only if $A$ is of projective type in the sense of Truss. In that case every positive finite power does. This gives an affirmative answer to Question 6.1 and a negative answer to Question 6.2 of Hu, Mao and Shen. The forward implication uses a countable language recording a single map; no enumeration of $A$ is required.
\end{abstract}
\section{Statement and coding}
An amorphous set is infinite with only finite and cofinite subsets. A set is dually Dedekind infinite if it has a noninjective self-surjection. Projective type means that $A$ admits a nondegenerate, locally homogeneous, locally finite pregeometry, as defined in [1, Section 3]. Degenerate means that the closure of every nonempty finite set is the union of the closures of its individual points.
\paragraph{Theorem (ZF).} For an amorphous $A$, the following are equivalent: (i) $A$ is of projective type; (ii) $\mathcal P(A)$ is dually Dedekind infinite; (iii) $\mathcal P(A)^k$ is dually Dedekind infinite for some positive finite $k$. In (iii) one may equivalently say every positive finite $k$.
\paragraph{Known direction.} Hu, Mao and Shen [1, Theorem 3.2] prove (i) implies (ii). A noninjective self-surjection acts on one coordinate of any positive power, giving the implication to every $k$. We prove (iii) implies (i).
Fix $k$ and a self-surjection $f$ as in (iii). Code each coordinate by a bit (finite or cofinite) and its finite part (itself or its complement). This representation is unique. Use one relation symbol for each pair of input/output bit vectors and each pair of lists of $k$ finite lengths. The relation on the concatenated tuple of entries says that the subtuples enumerate their respective finite parts without repetitions and that $f$ sends the coded input to the coded output. Repetitions across different coordinates are allowed. Nullary relations, or sentences coding them, handle empty lists. This gives a structure $M$ on $A$ in a specified countable language. Write $T=\mathrm{Th}(M)$.
\section{Finite types and homogeneous closure}
We give details of the model theory used on a possibly non-well-orderable domain; see [2] for the categoricity background. Every map from $A$ to a linear order has finite image. An infinite image would itself be amorphous by taking inverse images, whereas an infinite linearly ordered amorphous set is impossible: for each point either its set of predecessors or of successors is finite. On each side the respective finite cardinalities distinguish the points, giving an injection into $2\times\omega$, which contains no amorphous subset. By induction every map from $A^r$ into a linear order has finite image. Indeed each last-coordinate row has finite image; the row sizes are bounded by the induction hypothesis applied to their sizes on $A^{r-1}$. Enumerate each finite row image in increasing order, padding with a single fixed value if necessary, and apply induction to each of the boundedly many positions.
\paragraph{Isolated types.} The complete type of an $r$-tuple is a binary sequence in the lexicographic order on $2^\omega$. Therefore only finitely many such types occur. Separate each of them by a finite conjunction of distinguishing formulas. Their disjunction covers every tuple in $M$, hence in every model of $T$. These formulas isolate complete types and give finite-tuple homogeneity in the following precise sense: any partial elementary map on a finite tuple can be extended to include one more prescribed domain element, since the joint type has an isolating formula. The same argument works between two models of $T$.
\paragraph{Strong minimality.} For each unary formula with a parameter tuple, all its instances in $M$ are finite or cofinite. The sizes of their finite sides are uniformly bounded by the finite-image observation on the parameter tuples. The finite-or-cofinite assertion with that bound is a sentence of $T$. Thus $T$ is strongly minimal. Algebraic closure, denoted by $\mathrm{acl}$, is an exchange pregeometry in a strongly minimal theory. For clarity, if $b$ is algebraic over $S,a$ but not over $S$, witness this by a relation over $S$ with at most $m$ choices of $b$ for each $a$ (restrict its domain to such rows). There are at most $m$ parameters $b'$ whose inverse fibers are cofinite, since $m+1$ cofinite fibers intersect. The set of those parameters is definable over $S$ by the uniform finite-or-cofinite bound. Hence $b$ is not one of them, so $a$ lies in a finite fiber over $S,b$. This gives exchange. The remaining closure properties follow from finite definability of algebraicity.
For each fixed finite tuple $s$, there are finitely many $1$-types over $s$, by isolation of joint types. The algebraic ones have finitely many realizations, so $\mathrm{acl}(s)$ is finite. For each fixed arity the relation $x\in\mathrm{acl}(s)$ is uniformly definable by a union of finitely many joint types: finite fiber size is invariant under equal parameter types. These definitions agree in all models of $T$.
\paragraph{Local homogeneity in Truss's definition.} Let $S\subseteq D$ be finite closed sets and $a,b\in D\setminus S$. Every nonalgebraic point over a finite tuple has the unique generic type over it, determined by its cofinite definable sets. Extend $a,b$ separately to ordered bases of $D$ over $S$, of equal length by exchange. Matching these bases over the pointwise fixed $S$ is elementary, by successively matching generic types. Extend the match one point at a time to all of $D$ by finite-tuple homogeneity. Its image is in $D$, since $D$ is the closure of either basis over $S$; finiteness makes the image all of $D$. For any finite closed $U\supseteq D$, take an ordered basis of $U$ over $D$, extend the resulting map on $D$ by fixing that basis pointwise (again using generic types), and extend over its closure in the same way. The resulting permutations preserve closure on all subsets of the respective finite sets. This verifies the extension clause of local homogeneity without simultaneous choices. Thus a nondegenerate $\mathrm{acl}$ here already witnesses (i).
\section{The countable model in the degenerate case}
Suppose toward a contradiction that the pregeometry just obtained is degenerate. Its degeneracy is a scheme of first-order sentences, one for each tuple length by the uniform closure definitions. Write $C=\mathrm{acl}(\varnothing)$. Outside this finite set the sets $\mathrm{acl}(a)\setminus C$ partition the universe into blocks: membership between generic points is symmetric by exchange, and their singleton closures then coincide. All blocks have a common finite size $d$ by uniqueness of the generic type. Closure of a finite tuple is $C$ together with the blocks it meets.
\paragraph{A model with a full symmetric action.} There is a countable model $N$ of $T$ with its blocks enumerated in $d$ positions in such a way that \emph{every} permutation of the blocks, acting positionwise and fixing $C$, is an automorphism. Here are details that avoid choosing enumerations on $A$. Work with the real coding $T$ in the inner model $L[T]$, where choice holds. The isolation, finite-fiber bounds, and closure scheme above are properties of the first-order sentences in $T$ and hold there as well. One may name the finite tuple enumerating the core, using its isolated type. In a countable model take arbitrarily long lists of enumerated distinct blocks. Finite Ramsey's theorem and countable compactness give a model with an order-indiscernible sequence of such tuples indexed by $\mathbb Q$ over the core: require equality of truth values on increasing lists, and apply finite Ramsey to the finitely many formulas of each finite subtheory. In a stable theory every infinite order-indiscernible sequence on a dense order is totally indiscernible. Recall that an order-sensitive formula on a finite list yields, by an adjacent transposition and by fixing the other entries, a formula defining the ordering of a subinterval of the indiscernibles, the order property. A strongly minimal theory is stable: over a countable model there are just its algebraic $1$-types and its unique generic $1$-type, giving $\omega$-stability. These standard stability facts are applied inside $L[T]$ [3].
Take the union of the displayed blocks and the core. This set is algebraically closed by degeneracy and is infinite. It is an elementary substructure by Tarski--Vaught: each nonempty unary definable fiber with parameters in the union either is finite and contained in their algebraic closure, or is cofinite and meets the union. This gives $N$. Total indiscernibility means every permutation of the blocks acting in the stated way preserves every named predicate. The same remains true in the ambient ZF universe, including for permutations absent from $L[T]$, because preservation is tested on finite tuples.
\section{Transporting the map and bounding its orbits}
On $N$ interpret the named predicates on $k$-tuples of coded finite or cofinite sets, that is, on $D_N=(2\times[N]^{<\omega})^k$. This produces a total, surjective, noninjective map $g$. Here the unbounded list of possible output lengths requires justification. Given an input tuple in $N$, realize its isolated complete type in $M$. There $f$ has an output with some particular lengths and flags. Existence of that output is a single first-order existential formula with the input as parameters, so it holds for the input in $N$. For each fixed pair of shapes, uniqueness, including invariance under enumeration order, is first-order and transfers. The same argument fixing an output gives surjectivity. A particular collision of $f$ has finite codes of definite shapes, so its existential sentence transfers as well. This procedure never equates the full powerset of $N$ with coded finite and cofinite sets. By construction $g$ commutes with the positionwise symmetric action.
\paragraph{Uniform bound.} Fix $s\in D_N$. A block has one of $q=2^{dk}$ input patterns recording membership of each of its $d$ positions in the $k$ finite parts. The subgroup permuting blocks within each pattern fixes $s$ and hence every forward iterate $g^j(s)$. In any fixed output under this subgroup, the output pattern must be constant on each input pattern class, since individual blocks with the same input pattern can be interchanged. Its flags, its finite parts on the core, and at most $q$ output patterns determine it. Thus the orbit has at most
\[
K=2^{k(1+|C|)}q^q
\]
points, independently of $s$ or its finite-part sizes. For any map with this bound one has $g^{K+K!}=g^K$, by the pigeonhole principle applied on each orbit. Surjectivity of $g^K$ gives $g^{K!}=\mathrm{id}$, contradicting its noninjectivity. This excludes degeneracy, proving the theorem.
\section{The two original questions}
Question 6.1 is the case $k=1$ of the new implication. For Question 6.2, already for any single amorphous set $A$ and positive $n$, dual Dedekind infinity of $\mathcal P(A)^{n+1}$ forces projective type, whence dual Dedekind infinity of $\mathcal P(A)^n$. Thus the proposed family cannot exist. Both conclusions use the same theorem.
\paragraph{Sources.} [1] Y. Hu, R. Mao and G. Shen, \emph{Amorphous sets and dual Dedekind finiteness} (2025 preprint, 2026 journal), \url{https://arxiv.org/abs/2510.13508}, \href{https://doi.org/10.1016/j.apal.2026.103723}{journal version}. Their Theorem 3.2 gives one direction and Section 6 asks both questions. The original projective-type notion is due to J. K. Truss, \href{https://www.sciencedirect.com/science/article/pii/016800729400024W}{\emph{The structure of amorphous sets}} (1995).
[2] A. C. Walczak-Typke, \href{https://doi.org/10.2178/jsl/1129642121}{\emph{The first-order structure of weakly Dedekind-finite sets}} (2005).
[3] A. Pillay, \href{https://www3.nd.edu/~apillay/pdf/lecturenotes_modeltheory.pdf}{\emph{Lecture notes: Model theory}}; A. Chernikov, \href{https://www.math.ucla.edu/~chernikov/teaching/StabilityTheory285D/StabilityNotes.pdf}{\emph{Lecture notes on stability theory}}. The application of the countable block model to the self-surjection is the proposed converse. Independent proof and priority review is pending.
\end{document}
